% !TeX program = tectonic
\documentclass[11pt,a4paper]{article}
\usepackage{fontspec}
\usepackage[english]{babel}
\babelfont{rm}[Language=Default]{Liberation Serif}
\babelfont[Language=Default]{sf}{Carlito}
\babelfont[Language=Default]{tt}{DejaVu Sans Mono}
\usepackage{amsmath,amssymb,amsthm}
\usepackage{geometry}
\geometry{a4paper, top=2.4cm, bottom=2.4cm, left=2.4cm, right=2.4cm}
\usepackage{booktabs,tabularx,enumitem,xcolor,tcolorbox}
\tcbuselibrary{breakable,skins}
\definecolor{linkblue}{HTML}{1F4E79}
\definecolor{statgreen}{HTML}{2F6B3A}
\definecolor{goldacc}{HTML}{8B7E5A}
\usepackage[colorlinks=true,linkcolor=linkblue,urlcolor=linkblue,unicode=true]{hyperref}
\theoremstyle{plain}
\newtheorem{theorem}{Theorem}
\newtheorem{lemma}[theorem]{Lemma}
\newtheorem{proposition}[theorem]{Proposition}
\newtheorem{corollary}[theorem]{Corollary}
\theoremstyle{definition}
\newtheorem{definition}[theorem]{Definition}
\theoremstyle{remark}
\newtheorem{remark}[theorem]{Remark}
\newtcolorbox{statusbox}[1][]{colback=statgreen!5,colframe=statgreen!75,
  title={\textbf{Summary of results:} #1},fonttitle=\small,boxrule=0.7pt,arc=2pt,breakable}
\newtcolorbox{databox}[1][]{colback=goldacc!6,colframe=goldacc!80,
  title={\textbf{Protocol data:} #1},fonttitle=\small,boxrule=0.7pt,arc=2pt,breakable}
\newtcolorbox{verifybox}[1][]{colback=linkblue!5,colframe=linkblue!80,
  title={\textbf{Verification:} #1},fonttitle=\small,boxrule=0.7pt,arc=2pt,breakable}
\widowpenalty=10000
\clubpenalty=10000
\setlength{\parskip}{0.35em}
\newcommand{\CC}{\mathbb{C}}
\newcommand{\RR}{\mathbb{R}}
\newcommand{\ZZ}{\mathbb{Z}}
\newcommand{\QQ}{\mathbb{Q}}
\newcommand{\Ga}{\Gamma}
\newcommand{\Bfun}{\mathrm{B}}
\newcommand{\im}{\operatorname{Im}}
\newcommand{\re}{\operatorname{Re}}
\newcommand{\lcm}{\operatorname{lcm}}
\newcommand{\SNF}{\operatorname{SNF}}
\newcommand{\Jac}{\operatorname{Jac}}
\newcommand{\rank}{\operatorname{rank}}
\newcommand{\Ind}{\operatorname{Index}}
\newcommand{\disc}{\operatorname{disc}}
\begin{document}
\begin{center}
{\LARGE\bfseries Theorem 1: The Riemann Relations as an Identity}\\[0.4em]
{\large symmetry and positivity of the period matrix}\\[0.6em]
{\normalsize Isaev Iskhak Khamzatovich}\\[0.2em]
{\normalsize The <<Dynamic Principle>> Program --- the \texttt{hodge-laboratory}}\\[0.15em]
{\small Standalone theorem monograph · Монография 4/16}
\end{center}
\vspace{0.4em}
{\color{linkblue}\hrule height 0.8pt}
\vspace{0.8em}

\section{Setting}

The Riemann relations --- symmetry and positive definiteness of the
imaginary part of the period matrix --- are proved on a curve
\emph{identically}, with no numerical check: a numerical control of
an identity would be weaker than the proof. For the $N=15/30$ stand
this means: on the $91\times182$ and $406\times812$ matrices there
is nothing to check numerically --- it is proved.

\begin{definition}[normalized period matrix]
In the symplectic basis
$\mathfrak{B}=\{a_1,\dots,a_g,b_1,\dots,b_g\}$ normalize the forms by
$\int_{a_i}\omega_j=\delta_{ij}$ and set
$\Omega_{ij}=\int_{b_i}\omega_j$.
\end{definition}

\begin{theorem}[Riemann relations]\label{th:riemann}
(1) $\Omega^{\mathsf{T}}=\Omega$; (2) $\im\Omega$ is positive
definite: for every $0\ne\lambda\in\RR^g$,
$\lambda^{\mathsf{T}}(\im\Omega)\lambda>0$.
\end{theorem}

\begin{proof}
The tool is Riemann's bilinear identity: for the fundamental
$4g$-gon $F$ and closed forms $\alpha,\beta$,
\[
  \iint_F \alpha\wedge\beta
  =\sum_{i=1}^{g}\Bigl(
    \int_{a_i}\alpha\int_{b_i}\beta
    -\int_{b_i}\alpha\int_{a_i}\beta\Bigr).
\]
\emph{Symmetry:} at $\alpha=\omega_j$, $\beta=\omega_k$ the
holomorphic $(2,0)$-object $\omega_j\wedge\omega_k$ vanishes
($\Omega^2_{C_N}=0$), giving $\Omega_{jk}-\Omega_{kj}=0$.
\emph{Positivity:} at $\alpha=\omega$, $\beta=\overline{\omega}$ we
have $\omega\wedge\overline{\omega}=-2i|g|^2dx\wedge dy$; the left
side is $-2i\int|g|^2dxdy$; the right is
$-2i\lambda^T(\im\Omega)\overline{\lambda}^T$ (symmetry of part 1
used). Canceling $-2i$: the quadratic form equals
$\int_{C_N}|g|^2dxdy>0$.
\end{proof}

\begin{remark}[why no self-adjoint operator]
``Symplectic'' refers exclusively to the skew intersection form $J$
on the integer lattice $H_1(C_N,\ZZ)$ --- a topological object. The
theorem involves exactly two constructed objects: the cycle lattice
and the integrals of forms. The Riemann relations are an identity
that our own period matrix must satisfy; the numbers remain a
one-off sanity line outside the text.
\end{remark}

\begin{statusbox}[summary]
The theorem is proved by Riemann's bilinear identity; the only
analytic premise is Stokes' formula. Consequence: $J$ is a
polarization --- the input of the Hodge lattice index module
(theorem 3).
\end{statusbox}

\begin{verifybox}[reproduction]
The identity nature means: verification is a machine proof of
structure, not a numerical run. Lean checks the integer corollaries
(ranks $91$, $406$); the laboratory runs the sanity line (V7: rank
$=g$, condition numbers $8.63$, $18.17$).
\end{verifybox}

\vfill
{\color{linkblue}\hrule height 0.6pt}
\vspace{0.5em}
{\small\color{gray}
License: individual exclusive license of Isaev Iskhak Khamzatovich.
All rights to the computations, formulas, and texts belong to the author.
Verification: \texttt{python3 laboratory.py} (``Stands''/``Certificates''),
Lean: \texttt{verification/lean/HodgeLaboratory.lean}.
}
\end{document}
